\section*{Instrumentation Amplifier}
Si assumano gli operazionali ideali e le loro uscite
lineari. Si considerino inoltre i valori delle resistenze
pari a quelli nominali (e quindi R1=…=R4=R=10 kW).

\paragraph{Richiesta (a)}
Nell'ipotesi che il ramo di RG sia aperto, si calcoli il guadagno di modo differenziale
dell'amplificatore e se ne verifichi l'accordo con quanto ottenuto al punto 1b (suggerimento:
si consiglia di sfruttare la validità dei teoremi per le reti lineari e/o ricondurre il circuito ad
uno schema a blocchi di funzioni già note).

\paragraph{Richiesta (b)}
Si mostri che per ogni RG si abbia VOUT=0 se V1=V2, ovvero che il valore atteso per il
guadagno di modo comune Ac sia nullo.

\paragraph{Richiesta (c)}
Si calcoli più in generale l'espressione del guadagno di modo differenziale per valori finiti
di RG esplicitandone la dipendenza dal rapporto adimensionale $k=R/RG$.

\paragraph{Richiesta (d)}
Si calcoli quindi il valore atteso di Ad per RG=10, 5kW ($\implies k=1,2$) e lo si confronti con
quanto ottenuto al punto 2b.
